\section{Introduction}
\label{sec:intro}
\IEEEPARstart{I}{n} recent years, the paradigm of \textit{wifi movement detection} has transformed from an experimental curiosity into a foundational capability for ubiquitous ambient intelligence, smart healthcare monitoring, and non-intrusive surveillance. Unlike traditional vision-based perception systems that require line-of-sight illumination and introduce severe privacy and surveillance liabilities, wireless radio frequency (RF) signals natively permeate walls, preserve user anonymity, and operate continuously under all lighting conditions.

Despite these compelling advantages, achieving robust performance in practical environments remains a formidable open challenge \cite{2604.05475v1}. As demonstrated by recent pioneering studies \cite{1004.4588v1}, wireless propagation is governed by complex physical multipath reflections, temporal channel drift, and carrier frequency offsets. Large language and vision-language models have reached striking capabilities by training on internetscale data [2, 12]. A key enabling factor is the abundance of self-supervised training signal — text can be scraped at scale, and images can be paired with captions from alt-text and surrounding context. Yet for manybehavioralmodalities — eye movemen...

To overcome these critical barriers, in this paper we propose \textbf{RobustSense}, an end-to-end domain-invariant architecture tailored for \textit{wifi movement detection}. Specifically, we leverage physical-layer signal decoupling combined with self-supervised spatial contrastive learning. differ from comprehension; expertise is reflected in gaze allocation [13, 9]. Third, gaze is a fundamental signal across applications: training vision-language models on behavioral video, gaze-contingent rendering in virtual and augmented reality, social robotics, eye-based communication interfaces ...

The principal contributions of this work are summarized as follows:
\begin{itemize}
    \item \textbf{Physical-Layer Multipath Decoupling:} We formulate a rigorous mathematical preprocessing framework that cleans raw channel state indicators and separates static clutter from dynamic human motion Doppler shifts.
    \item \textbf{Domain-Invariant Feature Representation:} We introduce a self-supervised contrastive learning objective that extracts room-invariant spatial features, eliminating cumbersome offline recalibration when transitioning across distinct physical spaces.
    \item \textbf{Comprehensive Empirical Validation:} We benchmark our proposed framework against state-of-the-art architectures across multiple testbeds and public datasets, demonstrating marked improvements in cross-domain transfer accuracy and inference efficiency.
\end{itemize}

\textit{Organization:} The remainder of this paper is structured as follows. Section~\ref{sec:related_work} surveys the state of the art and prior literature. Section~\ref{sec:methodology} presents the formal system model and our proposed architecture. Section~\ref{sec:experiments} reports empirical benchmark results and ablation studies. Finally, Section~\ref{sec:conclusion} concludes the paper with future research directions.
